% Dyadic Fox--Wu--Tate extension for the fixed125 integrated manuscript.
\subsection{Mod-four Wu classes and exact twisted Fox--Bockstein depth}
\label{subsec:dyadic-fox-wu-tate}

The canonical orientation of the preceding subsection admits three further
finite interpretations: its reduction modulo $4$ is the characteristic
class of the mod-$2$ cup form; its full integral Fox row is the exact
obstruction to lifting $H^1$ with arbitrary rank-one twists; and its
vanishing row produces an integral top-cell trace for local Tate
duality.  These observations separate the \emph{ordinary} Bockstein
loss of information from the \emph{twisted} Kummer lifting property.

\begin{remark}[Disjointness at the even $q=2$ boundary]
\label{rem:dyadic-even-infinity-boundary}
One may formally write $f=\infty$ in the even-rank
two-generator-image normal form (iv) of
Theorem~\ref{thm:dyadic-relator-cases}.
But then $x_3^{2^\infty}=1$ and the resulting
relator is exactly the $f=\infty$ relation of
the even-rank cyclic-image family (iii),
namely $x_1^2[x_1,x_2]\cdots[x_{d-1},x_d]$.
Its orientation image is $\{\pm1\}$ and
has square-index $2$, not $4$.
The restriction to finite $f$ in family (iv)
is therefore a disjoint normalization and
excludes no Demu\v{s}kin isomorphism type.
\end{remark}

\begin{theorem}[The dyadic Fox--Wu/Bockstein identity]
\label{thm:dyadic-fox-wu-bockstein}
Let $P$ be a finite-rank infinite Demu\v{s}kin pro-$2$ group
and let $\theta=\theta_P:P\to\mathbb Z_2^\times$
be its canonical dualizing orientation.  Define
\begin{equation}\label{eq:dyadic-fox-wu-class}
w_P:P\longrightarrow\mathbf F_2,\qquad
w_P(g)=\frac{\theta(g)-1}{2}\bmod2.
\end{equation}
This is a continuous group homomorphism, hence a class in
$H^1(P,\mathbf F_2)$.  For every
$\chi\in H^1(P,\mathbf F_2)$ one has
\begin{equation}\label{eq:dyadic-fox-wu-identity}
\boxed{\qquad
\beta_1(\chi)=\chi\smile\chi=\chi\smile w_P
\quad\text{in }H^2(P,\mathbf F_2).\qquad}
\end{equation}
In the normal-form basis of
Theorem~\ref{thm:dyadic-relator-cases},
\begin{equation}\label{eq:dyadic-wu-basis-values}
w_P=
\begin{cases}
0,&q\ne2,\\
\chi_1,&q=2,\ d\text{ odd (case (ii))},\\
\chi_2,&q=2,\ d\text{ even (cases (iii),(iv))}.
\end{cases}
\end{equation}
For a finite extension $K/\mathbf Q_2$ and
$P=G_K(2)$, the class $w_P$ equals the
Kummer character $\kappa_{-1}$ of $-1\in K^\times$:
\begin{equation}\label{eq:dyadic-kummer-minus-one}
\chi\smile\chi=\chi\smile\kappa_{-1}.
\end{equation}
In particular, $w_P=0$ precisely when $\mu_4\subset K$
in the local Galois case.
\end{theorem}

\begin{proof}
Every $2$-adic unit is $1$ modulo $2$, so the formula
\eqref{eq:dyadic-fox-wu-class} is defined.  Multiplication
of two elements $1+2a$ and $1+2b$ gives
$1+2(a+b)+4ab$, proving additivity modulo $2$.
The ordinary connecting Bockstein for
$0\to\mathbf F_2\to\mathbf Z/4\to\mathbf F_2\to0$
equals the degree-one cup square:
$\beta_1(\chi)=\chi\smile\chi$.
Theorem~\ref{thm:dyadic-twisted-fox-orientation}
gives the reduction of $\theta$ modulo $4$:
it is trivial in case (i), is $-1$ on $x_1$ in
case (ii), and is $-1$ on $x_2$ in
cases (iii),(iv), with all other generator values
equal to $1$ modulo $4$.
This proves \eqref{eq:dyadic-wu-basis-values}.

The quadratic Fox cup matrices
\eqref{eq:dyadic-quadratic-symbols}
have zero diagonal in case (i), and otherwise
their only nonzero diagonal entry is $B_{11}=1$.
In case (ii), the first generator is orthogonal
to every other basis vector, so
$B(\chi,\chi)=B(\chi,\chi_1)$ for every $\chi$.
In cases (iii),(iv), the first block is
$\bigl(\begin{smallmatrix}1&1\\1&0\end{smallmatrix}\bigr)$,
so $B(\chi,\chi)=B(\chi,\chi_2)$.
This proves \eqref{eq:dyadic-fox-wu-identity}
basis-independently by bilinearity.

For $P=G_K(2)$ the orientation is the restricted
$2$-adic cyclotomic character.  Reducing it modulo
$4$ records the action of $G_K$ on a square root
$i$ of $-1$, namely the quadratic extension
$K(i)/K$.  Its mod-$2$ character is precisely
the Kummer class $\kappa_{-1}$.  Hence
$w_P=\kappa_{-1}$ and
\eqref{eq:dyadic-kummer-minus-one} follows.
The character is zero exactly when $i\in K$.
\end{proof}

For an arbitrary continuous unit-valued character
$\psi:P\to\mathbb Z_2^\times$, retain a fixed
minimal one-relator presentation
$P=\langle x_1,\ldots,x_d\mid r\rangle_{\mathrm{pro}-2}$
and define the \emph{twisted top-cell Fox row}
\begin{equation}\label{eq:dyadic-psi-fox-row}
J_\psi:=
\bigl(\varepsilon_\psi(\partial_1r),\ldots,
\varepsilon_\psi(\partial_dr)\bigr)\in\mathbb Z_2^d,
\qquad
a(\psi):=\min_i v_2((J_\psi)_i),
\end{equation}
where $v_2(0)=\infty$.  When $a(\psi)<\infty$,
write
\begin{equation}\label{eq:dyadic-psi-leading-row}
j_\psi=
2^{-a(\psi)}J_\psi\bmod2
\in(\mathbf F_2)^d\setminus\{0\}.
\end{equation}
The ideal generated by $J_\psi$ is invariant under
Fox presentation change; individual row coordinates
and the hyperplane below transform covariantly.

\begin{theorem}[Exact twisted Fox lifting-depth obstruction]
\label{thm:dyadic-twisted-fox-lifting-depth}
Let $P$ and $\psi$ be as above.  Let
$\mathbb Z_2(\psi)$ be the rank-one integral module
on which $P$ acts through $\psi$ and put
\begin{equation}\label{eq:dyadic-psi-lifting-def}
L_m^\psi(P)=
\operatorname{im}\left(
H^1(P,\mathbb Z_2(\psi)/2^m)
\longrightarrow H^1(P,\mathbf F_2)\right),
\quad m\ge1.
\end{equation}
Then the rank-one twisted Fox cochain complex is
\begin{equation}\label{eq:dyadic-psi-three-term-fox}
\mathbb Z_2
\xrightarrow{\ \bigl(\psi(x_i)-1\bigr)_i\ }
\mathbb Z_2^d
\xrightarrow{\ J_\psi\ }
\mathbb Z_2.
\end{equation}
Consequently
\begin{equation}\label{eq:dyadic-psi-h2}
H^2(P,\mathbb Z_2(\psi))
\cong
\begin{cases}
\mathbb Z_2,&a(\psi)=\infty,\\
\mathbb Z_2/2^{a(\psi)}\mathbb Z_2,
  &a(\psi)<\infty,
\end{cases}
\end{equation}
and the entire lifting filtration is
\begin{equation}\label{eq:dyadic-psi-lifting-formula}
L_m^\psi(P)=
\begin{cases}
H^1(P,\mathbf F_2),
 &m\le a(\psi),\\
\ker\bigl(j_\psi:(\mathbf F_2)^d\to\mathbf F_2\bigr),
 &m>a(\psi),\quad a(\psi)<\infty.
\end{cases}
\end{equation}
Here the identification $H^1(P,\mathbf F_2)
\cong(\mathbf F_2)^d$ uses the minimal generators.
For all $\psi$, either $a(\psi)=\infty$ or
$a(\psi)\ge1$.  Moreover, the following conditions
are equivalent:
\begin{equation}\label{eq:dyadic-psi-orientation-criterion}
\boxed{\quad
\psi=\theta_P
\quad\Longleftrightarrow\quad
J_\psi=0
\quad\Longleftrightarrow\quad
H^2(P,\mathbb Z_2(\psi))\simeq\mathbb Z_2
\quad\Longleftrightarrow\quad
L_m^\psi(P)=H^1(P,\mathbf F_2)\ \forall m.
\quad}
\end{equation}
Thus the canonical orientation is also the unique
integral unit twist with \emph{unrestricted all-level
mod-$2$ degree-one liftability}.
\end{theorem}

\begin{proof}
Apply the completed finite free Fox resolution
\eqref{eq:fox-chain-complex} to the rank-one module
$\mathbb Z_2(\psi)$.  The displayed differentials
follow from the Fox cochain formulas
\eqref{eq:fox-d0}--\eqref{eq:fox-d1}.
Since the top cohomology is the cokernel of
the single row $J_\psi$, its Smith invariant
over the discrete valuation ring $\mathbb Z_2$
gives \eqref{eq:dyadic-psi-h2}.

The presentation is minimal and every unit
$\psi(x_i)$ is $1$ modulo $2$, so both Fox
differentials vanish modulo $2$.  In particular,
$H^1(P,\mathbf F_2)\cong(\mathbf F_2)^d$ and
$J_\psi\equiv0\bmod2$.
If $m\le a(\psi)$, the twisted top differential
vanishes modulo $2^m$, so every mod-$2$ degree-one
cocycle lifts.  If $m>a(\psi)$, reducing the
equation $J_\psi v=0\bmod2^m$ after dividing by
$2^{a(\psi)}$ gives $j_\psi\bar v=0$.
Conversely, if $j_\psi\bar v=0$, choose an index
where $2^{-a(\psi)}(J_\psi)_i$ is a unit.  Keeping
all other entries fixed and adjusting that entry
modulo $2^{m-a(\psi)}$ solves the lifted equation
without changing $\bar v$.  Hence the image is
exactly the stated hyperplane.

Theorem~\ref{thm:dyadic-twisted-fox-orientation}
identifies $\theta_P$ as the \emph{unique} unit
character annihilating the entire twisted Fox row.
This proves the first two equivalences.
Equation~\eqref{eq:dyadic-psi-h2} proves the third,
and \eqref{eq:dyadic-psi-lifting-formula} the fourth:
if $a(\psi)<\infty$, its leading row is nonzero
and full liftability fails first at $m=a(\psi)+1$.
\end{proof}

\begin{corollary}[The orientation repairs every untwisted Bockstein loss]
\label{cor:dyadic-orientation-kummer-lifts}
For every such $P$ and every $m\ge1$,
\[
H^1(P,\mathbb Z_2(\theta_P)/2^m)
\twoheadrightarrow H^1(P,\mathbf F_2).
\]
If the trivial twist has invariant
$\mathfrak q(P)=(2^f)$ with finite $f$,
then its lifting filtration loses exactly
one mod-$2$ direction at level $f+1$,
whereas the canonical twist loses none.
For the local Galois groups $G_K(2)$,
the twisted surjectivity is the finite-Fox
realization of the Kummer lifting property.
\end{corollary}

\begin{proof}
For the canonical orientation the top Fox row
vanishes integrally by
Theorem~\ref{thm:dyadic-twisted-fox-orientation};
apply \eqref{eq:dyadic-psi-lifting-formula}.
For the trivial twist $J_1$ is the integral
exponent-sum row, whose ideal is $(2^f)$;
its leading reduction is nonzero, giving
the claimed codimension-one loss.
\end{proof}

\begin{theorem}[Integral dyadic Kummer--Tate groups from one Fox row]
\label{thm:dyadic-fox-kummer-tate}
Let $K/\mathbf Q_2$ be finite of degree $n$,
$P_K=G_K(2)$, and put
\[
q_K=|\mu_{2^\infty}(K)|=2^{f_K}\ge2,
\qquad d=n+2.
\]
The cyclotomic orientation $\theta_K$ of $P_K$
has both Fox rows
\begin{equation}\label{eq:dyadic-kummer-fox-rows}
J_{\theta_K}=0,\qquad
\bigl(\theta_K(x_i)-1\bigr)_i
\text{ generating the ideal }(q_K)\subset\mathbb Z_2.
\end{equation}
Thus its Fox cochain complex with cyclotomic
coefficients computes
\begin{equation}\label{eq:dyadic-kummer-tate-groups}
\boxed{
\begin{aligned}
H^0(G_K,\mathbb Z_2(1))&=0,\\
H^1(G_K,\mathbb Z_2(1))
  &\cong\mathbb Z_2^{\,n+1}
     \oplus\mathbb Z_2/q_K\mathbb Z_2,\\
H^2(G_K,\mathbb Z_2(1))&\cong\mathbb Z_2.
\end{aligned}}
\end{equation}
The untwisted Fox complex likewise gives
\begin{equation}\label{eq:dyadic-untwisted-tate-groups}
H^0(G_K,\mathbb Z_2)=\mathbb Z_2,\qquad
H^1(G_K,\mathbb Z_2)\cong\mathbb Z_2^{\,n+1},\qquad
H^2(G_K,\mathbb Z_2)\cong
   \mathbb Z_2/q_K\mathbb Z_2.
\end{equation}
In the $K=\mathbf Q_2$ basis of
Proposition~\ref{prop:dyadic-q2-explicit-anchor},
the entire twisted Fox differential is
\begin{equation}\label{eq:dyadic-q2-twisted-complex}
\mathbb Z_2
\xrightarrow{\;(-2,\ 0,\ -4/3)^{\rm t}\;}
\mathbb Z_2^3
\xrightarrow{\;(0,\ 0,\ 0)\;}
\mathbb Z_2.
\end{equation}
In particular,
$H^1(G_{\mathbf Q_2},\mathbb Z_2(1))
\cong\mathbb Z_2^2\oplus\mathbb Z/2$
and the top twisted cohomology is torsion-free,
whereas
$H^2(G_{\mathbf Q_2},\mathbb Z_2)
\cong\mathbb Z/2$.
\end{theorem}

\begin{proof}
The maximal pro-$2$ local Galois quotient
is Demu\v{s}kin, and its canonical orientation
is the cyclotomic character
\cite{LabuteClassification,NSW}.
Theorem~\ref{thm:dyadic-twisted-fox-orientation}
gives $J_{\theta_K}=0$.
By definition of the cyclotomic action,
$\theta_K(G_K)$ acts trivially on all $q_K$-th
roots of unity but not on all $2q_K$-th
roots, since $q_K$ is the exact order
of the $2$-primary roots of unity in $K$.
The ideal generated by the differences
$\theta_K(g)-1$ is therefore exactly $(q_K)$.
As the chosen $x_i$ topologically generate
$P_K$, the same ideal is generated by
$\theta_K(x_i)-1$.
Consequently the twisted Fox complex has
an injective first differential whose image
has one Smith invariant $q_K$, and a zero
second differential.  Its cohomology is
\eqref{eq:dyadic-kummer-tate-groups}.
This agrees with local Kummer theory,
$H^1(G_K,\mathbb Z_2(1))
\simeq\widehat{K^\times}^{\,2}$.

For the trivial twist, the first differential
is zero, and the second differential is
the exponent-sum row with ideal $(q_K)$,
by local class field theory and the
Demu\v{s}kin abelianization invariant.
Its kernel is free of rank $d-1$ and its
cokernel is $\mathbb Z_2/q_K$, proving
\eqref{eq:dyadic-untwisted-tate-groups}.
Finally the known $G_{\mathbf Q_2}(2)$
orientation takes values $(-1,1,-1/3)$
on $(a,s,y)$; subtracting $1$ gives
$(-2,0,-4/3)$, and its twisted Fox row
vanishes.  Formula
\eqref{eq:dyadic-q2-twisted-complex}
and the claimed groups follow by Smith
normal form.
\end{proof}

\begin{corollary}[Exact cyclotomic-power depth at $\mathbf Q_2$]
\label{cor:dyadic-q2-cyclotomic-power-depth}
Retain the marked generators $(a,s,y)$ of
Proposition~\ref{prop:dyadic-q2-explicit-anchor}.
For any odd integer $k\ge1$, let
$\psi_k=\theta_{\mathbf Q_2}^{\,k}$.  Its entire
twisted top Fox row is
\begin{equation}\label{eq:dyadic-q2-cyclotomic-power-row}
J_{\psi_k}=(0,\ 3-3^k,\ 0).
\end{equation}
For $k=1$ this row is zero.  For every odd
$k\ge3$ its first nonvanishing depth is
\[
a(\psi_k)=2+v_2(k-1),
\]
and therefore
\begin{equation}\label{eq:dyadic-q2-cyclotomic-power-lifting}
L_m^{\psi_k}(P_{\mathbf Q_2})=
\begin{cases}
H^1(P_{\mathbf Q_2},\mathbf F_2),
 &m\le 2+v_2(k-1),\\
\operatorname{span}_{\mathbf F_2}
   \{\chi_a,\chi_y\},
 &m>2+v_2(k-1).
\end{cases}
\end{equation}
In particular
\[
H^2(G_{\mathbf Q_2},\mathbb Z_2(k))
\cong \mathbb Z_2/2^{\,2+v_2(k-1)}\mathbb Z_2
\qquad(k\ge3\text{ odd}),
\]
while $H^2(G_{\mathbf Q_2},\mathbb Z_2(1))
\cong\mathbb Z_2$.
The exponent $2+v_2(k-1)$ is an exact
orientation-sensitive $2$-adic lifting depth,
not a new invariant of the untwisted group.
\end{corollary}

\begin{proof}
For odd $k$ one has
$\psi_k(a)=-1$, $\psi_k(s)=1$,
and $\psi_k(y)=-3^{-k}$.
The first Fox derivative of $a^2$
evaluates to $1+\psi_k(a)=0$.
The $y$-derivative of $s^4[s,y]$
vanishes because $\psi_k(s)=1$.
Finally the $s$-derivative evaluates to
$4+\psi_k(y)^{-1}-1=3-3^k$.
This proves \eqref{eq:dyadic-q2-cyclotomic-power-row}.
For even $n=k-1\ge2$, the $2$-adic
lifting-the-exponent formula gives
$v_2(3^n-1)=2+v_2(n)$.
The leading Fox row therefore reduces to
the $\chi_s$ coordinate, and
Theorem~\ref{thm:dyadic-twisted-fox-lifting-depth}
gives the precise filtration and top
cohomology claimed.
\end{proof}
\begin{proposition}[Transfer-normalized integral Fox--Tate pairing]
\label{prop:dyadic-sylow-fox-tate-pairing}
Let $E/\mathbf Q_2$ be finite, let
$T$ be a finite free $\mathcal O_E$-lattice
with continuous $G_K$-action, and put
$T^\vee(1)=\operatorname{Hom}_{\mathcal O_E}
(T,\mathcal O_E)(1)$.
Choose a residual Sylow-$2$ preimage
$H\subseteq G_K$ of odd index $m$, and let
$K_H$ be its fixed field.
The orientation-twisted top Fox complex of
$P_H=H(2)$ has zero final differential and
a top-degree trace
\[
\operatorname{tr}_{H}^{\rm Fox}:
C^2_{\rm Fox}(P_H,\mathcal O_E(1))
\longrightarrow\mathcal O_E
\]
normalized to the local invariant
$H^2(G_{K_H},\mathcal O_E(1))
\simeq\mathcal O_E$.
With any Fox diagonal approximation and the
coefficient evaluation
$T\otimes T^\vee(1)\to\mathcal O_E(1)$,
the formula
\begin{equation}\label{eq:dyadic-transfer-tate-pairing}
\boxed{\quad
\langle a,b\rangle_K^{\rm Fox}
=
\frac1m\operatorname{tr}_{H}^{\rm Fox}
\left(
\operatorname{res}_{H}(a)
  \smile_{\rm Fox}
\operatorname{res}_{H}(b)
\right)\quad}
\end{equation}
represents the intrinsic integral local
Tate pairing of perfect complexes
\[
R\Gamma(G_K,T)\otimes_{\mathcal O_E}^{\mathbf L}
R\Gamma(G_K,T^\vee(1))
\longrightarrow\mathcal O_E[-2].
\]
In particular it is perfect in the derived
sense, independent of the Sylow subgroup and
presentation, and requires no integral
termwise $\Delta_K$-invariants.
Its finite Fox matrix representative depends
on the chosen diagonal approximation and
normalization; the derived pairing does not.
\end{proposition}

\begin{proof}
Theorem~\ref{thm:dyadic-fox-kummer-tate},
applied to $K_H$ and extended to $\mathcal O_E$,
gives $H^2(P_H,\mathcal O_E(1))\cong\mathcal O_E$
and shows that the twisted top Fox differential
is zero.  Choose its generator so that the
resulting trace coincides with the usual
local invariant under the Fox--Galois comparison.
The finite Fox diagonal and tensor Fox
formula of
Theorem~\ref{thm:twisted-tensor-fox-cup}
represent the intrinsic cup product.
For a finite extension $K_H/K$ of degree
$m$, compatibility of local invariants
with restriction gives
\[
\operatorname{inv}_{K_H}(\operatorname{res}_{H}\eta)
=m\,\operatorname{inv}_{K}(\eta)
\quad
(\eta\in H^2(G_K,\mathcal O_E(1))).
\]
Since $m$ is odd, division by $m$ is integral.
Hence \eqref{eq:dyadic-transfer-tate-pairing}
is the original cup/evaluation/invariant
pairing on $G_K$.
Integral local Tate duality identifies its
adjoint with a derived equivalence
\[
R\Gamma(G_K,T^\vee(1))
\simeq
R\operatorname{Hom}_{\mathcal O_E}
(R\Gamma(G_K,T),\mathcal O_E)[-2]
\]
\cite{NSW}.
Different Sylow preimages and Fox diagonals
give representatives of this same intrinsic
pairing.  In particular, no assertion that
normalized corestriction is itself a
multiplicative cochain map is used.
\end{proof}
